\documentclass[10pt,a4paper]{amsart}
\usepackage[margin=19mm]{geometry}
\usepackage{amsmath,amssymb,mathtools}
\usepackage[scr=rsfs,cal=euler]{mathalfa}
\usepackage{xfrac}
\usepackage[unicode,hidelinks,bookmarks=false]{hyperref}
\newcommand{\ie}{i.e.\ }
\newcommand{\cf}{cf.\ }
\newcommand{\resp}{respectively\ }
\newcommand{\Subsection}{\subsection}
\providecommand{\nobreakdash}{\nobreak\hbox{-}}
\setlength{\emergencystretch}{2em}
\multlinegap0pt
\usepackage{amssymb}
\usepackage[shortlabels]{enumitem}
\usepackage{tikz}
\usepackage{float}
\newtheorem{prop}{Proposition}
\title{Asymptotic consensus and flocking under decaying persistent excitation on rooted digraphs}
\author{Chiara Cicolani, Elisa Continelli, Cristina Pignotti}
\date{Source: arXiv:2608.27263v1 (CC BY 4.0)}
\begin{document}
\maketitle

\noindent\emph{Source and licence.} Asymptotic consensus and flocking under decaying persistent excitation on rooted digraphs. Version arXiv:2608.27263v1 (CC BY 4.0), distributed by arXiv under the Creative Commons Attribution 4.0 International licence, http://creativecommons.org/licenses/by/4.0/, as recorded in the arXiv licence metadata for this exact version and shown on the abstract page https://arxiv.org/abs/2608.27263v1. Full licence notice: \texttt{licenses/arxiv-2608.27263-CC-BY-4.0.txt}.\par\smallskip

\noindent\textit{Editorial scope.} Counted unit: the article's prop lemma3onoff (source0.tex lines 371--380). The article's complete original proof of it is delivered with it (source0.tex lines 384--493, one \texttt{proof} environment of 110 lines, which the article places away from the statement). No mathematics is written, repaired or re-derived: the statement and the proof are the article's own lines, in the article's own notation, and no retained line is substituted or re-flowed.  Retained uncounted as the source-local context the proof needs: lines 207--209; lines 223--225; lines 310--312; lines 313--315; lines 316--318; lines 334--336; lines 363--365.  Nothing was omitted from any retained span, so the package carries no omission record.  No retained line contains a citation, so the wrapper carries no bibliography and no bibliography entry.  No retained line contains a figure environment, an image-inclusion command or drawing code, so no image is delivered and the package declares no asset dependency.  Editorial formatting only, in addition: an AMS \texttt{amsart} preamble that restates the article's own declarations and macros used by the retained lines, namely the environments \texttt{prop} on separate counters, together with the standard packages those lines need.  The retained environments print in this document as Proposition 1 in order of appearance, whereas the article prints the same items with the numbering of its own full document.  The complete native source of the article as distributed in the arXiv e-print for this exact version is bundled unchanged as \texttt{sources/208701/source0.tex}.
	\begin{equation}\label{gamma}
		\gamma_r:=\max_{i =1,\dots, N} \text{dist}(r,i). 
	\end{equation}

	\begin{equation}\label{PE}
		\int_{t}^{t+T}\alpha_{ij}(s)ds\geq \frac{\tilde{\alpha}}{(1+t)^\beta},\quad \forall t\geq 0.
	\end{equation}

		\begin{equation}\label{In}
			I_n:=[n(\gamma_r (T+\tau)+\tau)-\tau,n(\gamma_r (T+\tau)+\tau)],
		\end{equation}

		\begin{equation}\label{m_n}
			m_n^v:=\min_{i=1,\dots,N}\min_{s\in I_n}\,\langle x_{i}(s),v\rangle,
		\end{equation}

		\begin{equation}\label{M_n}
			M_n^v:=\max_{j=1,\dots,N}\max_{s\in I_n}\,\langle x_{j}(s),v\rangle.
		\end{equation}

		\begin{equation}\label{scalpronoff}
			m_n^v\leq \langle x_{i}(t),v\rangle \leq M_n^v,
		\end{equation}for all  $t\geq n(\gamma_r(T+\tau)+\tau)-\tau$ and for any $i=1,\dots,N$, where $m_n^v$ and $M_n^v$ are defined in \eqref{m_n} and \eqref{M_n}, respectively, $\gamma_r$ is the distance from the root defined in \eqref{gamma} and $T$ is the positive constant in \eqref{PE}.

        In particular,        \begin{equation}\label{stima_psionoff}
			\psi (x_i(t), x_j(t-\tau_{ij}(t)))\ge \psi_{0}:=\min_{\vert y\vert, \vert z\vert \le  C_{0}}\psi(y,z)>0,
		\end{equation}

	\begin{prop}\label{lemma3onoff}
		For all $n\in \mathbb{N}_0$ and $v\in \mathbb{R}^d$, it holds
		\begin{equation} \label{Bonoff}
			m_{n}^v+\Gamma_n(\langle x_r(n(\gamma_r(T+\tau)+\tau)),v \rangle-m^v_{n}) \leq \langle x_i(t),v\rangle\leq M^v_{n}-\Gamma_n(M^v_{n}-\langle x_r(n(\gamma_r(T+\tau)+\tau)),v \rangle), 
		\end{equation}
		for all $t\in I_{n+1}$ and $i =1,\dots, N$, where $I_{n+1}$ is the time interval defined in \eqref{In} and		\begin{equation}\label{Gammaonoff}
			\Gamma_n:=e^{-K (\frac{1}{2}(\gamma_{r}^2+3\gamma_r)(T+\tau)+\tau)}\left(\frac{\psi_{0}\tilde{\alpha}}{(N-1)(1+n(\gamma_r(T+\tau)+\tau)+(\gamma_r-1)T+\gamma_r\tau)^\beta}\right)^{\gamma_r},
		\end{equation}
		being $m_n^v$ and $M_n^v$ the quantities defined in \eqref{m_n} and \eqref{M_n}, respectively, $\gamma_r$ the distance from the root defined in \eqref{gamma}, $\beta$, $T$ and $\tilde{\alpha}$ the constants in \eqref{PE}.
	\end{prop}

	\begin{proof}[Proof of Proposition \ref{lemma3onoff}]
		 Let $n\in \mathbb{N}_0$ and $v \in \mathbb{R}^d$. Let us fix an index $i=1,\dots,N$ with $i\neq r$. Due to the fact that the direct graph $\mathcal{G}$ is rooted, we can always find a path from $r$ to $i$. Let $i_0=r,i_1,\dots,i_p=i$, $1\leq p\leq \gamma_r$, be the shortest from $r$ to $i$. Then, for a.e. $t\in [n(\gamma_r(T+\tau)+\tau),(n+1)(\gamma_r(T+\tau)+\tau)]$, from \eqref{scalpronoff} we can write
		$$\begin{array}{l}
			\vspace{0.3cm}\displaystyle{\frac{d}{dt}\langle x_r(t),v\rangle=\sum_{j:j\neq r}\chi_{rj}(t)b_{rj}(t)(\langle x_j(t-\tau_{rj}(t)),v\rangle-\langle x_r(t),v\rangle)}\\
			\vspace{0.3cm}\displaystyle{\hspace{1.5cm}\leq \sum_{j:j\neq r}\chi_{rj}b_{rj}(t)(M_n^v-\langle x_r(t),v\rangle)}\\
			\displaystyle{\hspace{1.5cm}\leq \frac{K}{N-1}\sum_{j:j\neq r}(M_n^v-\langle x_r(t),v\rangle)=K(M_n^v-\langle x_r(t),v\rangle).}
		\end{array}$$
		The Gronwall's inequality yields
		\begin{equation}\label{primopasso}\begin{array}{l}
			\vspace{0.3cm}\displaystyle{\langle x_r(t),v\rangle\leq e^{-K(t-n(\gamma_r(T+\tau)+\tau))}\langle x_r(n(\gamma_r(T+\tau)+\tau)),v\rangle+M_n^v(1-e^{-K(t-n(\gamma_r(T+\tau)+\tau)})}\\
			\vspace{0.3cm}\displaystyle{\hspace{1.5cm}=M_n^v-e^{-K(t-n(\gamma_r(T+\tau)+\tau))}(M_n^v-\langle x_r(n(\gamma_r(T+\tau)+\tau)),v \rangle )}\\
			\displaystyle{\hspace{1.5cm}\leq M_n^v-e^{-K(\gamma_r(T+\tau)+\tau)}(M_n^v-\langle x_r(n(\gamma_r(T+\tau)+\tau)),v \rangle),}
		\end{array}
        \end{equation}
        for all $t\in [n(\gamma_r(T+\tau)+\tau),(n+1)(\gamma_r(T+\tau)+\tau)]$.
		\\Now, for a.e. $t \in [n(\gamma_r(T+\tau)+\tau)+\tau,(n+1)(\gamma_r(T+\tau)+\tau)]$, using \eqref{scalpronoff} and \eqref{primopasso} we get
		$$\begin{array}{l}
			\vspace{0.3cm}\displaystyle{\frac{d}{d t}  \langle x_{i_1}(t),v\rangle =  \sum_{j \neq i_1, r} \chi_{i_1j} b_{i_1j}(t)( \langle x_{j}(t-\tau_{i_1j}(t)),v\rangle-\langle x_{i_1}(t),v\rangle)+ b_{i_1r}(t)(\langle x_{r}(t-\tau_{i_1r}(t)),v\rangle-\langle x_{i_1}(t),v\rangle) }\\
			\vspace{0.3cm}\displaystyle{\hspace{1cm}\leq \sum_{j \neq i_1 ,r} \chi_{i_1j} b_{i_1j}(t)( M_{n}^v-\langle x_{i_1}(t),v\rangle)}\\
			\vspace{0.3cm}\displaystyle{\hspace{3cm}+b_{i_1r}(t)\left(M_n^v-e^{-K(\gamma_r(T+\tau)+\tau)}(M_n^v-\langle x_r(n(\gamma_r(T+\tau)+\tau)),v \rangle)-\langle x_{i_1}(t),v\rangle\right)}\\
			\vspace{0.3cm}\displaystyle{\hspace{1cm}=( M_{n}^v-\langle x_{i_1}(t),v\rangle)\sum_{j \neq i_1 , r} \chi_{i_1j} b_{i_1j}(t)}\\
			\displaystyle{\hspace{2cm}+ b_{i_1 r}(t)\left(M_n^v-e^{-K(\gamma_r(T+\tau)+\tau)}(M_n^v-\langle x_r(n(\gamma_r(T+\tau)+\tau)),v \rangle)-\langle x_{i_1}(t),v\rangle\right).}
		\end{array}$$
		Noticing that
		$$\sum_{j \neq i_1 ,r} \chi_{i_1j} b_{i_1j}(t)=\sum_{\substack{j \neq i_1 }} \chi_{i_1j} b_{i_1j}(t)-b_{i_1r}(t)\leq \frac{KN_{i_1}}{N-1}-b_{i_1r}(t),$$
		from \eqref{stima_psionoff} we obtain
		$$\begin{array}{l}
			\vspace{0.3cm}\displaystyle{\frac{d}{d t}  \langle x_{i_1}(t),v\rangle \leq \frac{KN_{i_{1}}}{N-1}(M_{n}^v-\langle x_{i_1}(t),v\rangle) -b_{i_1r}(t)(M_{n}^v-\langle x_{i_1}(t),v\rangle) }\\
			\vspace{0.3cm}\displaystyle{\hspace{3cm}+ b_{i_1r}(t)\left(M_n^v-e^{-K(\gamma_r(T+\tau)+\tau)}(M_n^v-\langle x_r(n(\gamma_r(T+\tau)+\tau)),v \rangle)-\langle x_{i_1}(t),v\rangle\right)}\\
			\vspace{0.3cm}\displaystyle{\hspace{1cm}=\frac{KN_{i_{1}}}{N-1}(M_{n}^v-\langle x_{i_1}(t),v\rangle) -e^{-K(\gamma_r(T+\tau)+\tau)}(M_n^v-\langle x_r(n(\gamma_r(T+\tau)+\tau)),v \rangle)b_{i_1r}(t)}\\
			\vspace{0.3cm}\displaystyle{\hspace{1cm}\leq \frac{KN_{i_{1}}}{N-1}(M_{n}^v-\langle x_{i_1}(t),v\rangle) -e^{-K(\gamma_r(T+\tau)+\tau)}(M_n^v-\langle x_r(n(\gamma_r(T+\tau)+\tau)),v \rangle)\alpha_{i_1r}(t)\frac{\psi_{0}}{N-1}}\\
			\displaystyle{\hspace{1cm}=\frac{KN_{i_{1}}}{N-1} M_n^v-e^{-K(\gamma_r(T+\tau)+\tau)}(M_n^v-\langle x_r(n(\gamma_r(T+\tau)+\tau)),v \rangle)\alpha_{i_1r}(t)\frac{\psi_{0}}{N-1}-\frac{KN_{i_{1}}}{N-1}\langle x_{i_1}(t),v\rangle.}
		\end{array}$$
		Hence, the Gronwall's inequality gives
		$$ \begin{array}{l}
			\vspace{0.3cm}\displaystyle{\langle x_{i_1}(t),v\rangle\leq e^{-\frac{KN_{i_1}}{N-1}(t-n(\gamma_r(T+\tau)+\tau)-\tau)} \langle x_{i_1}(n(\gamma_r(T+\tau)+\tau)+\tau),v\rangle+M^v_n(1-e^{-\frac{KN_{i_1}}{N-1}(t-n(\gamma_r(T+\tau)+\tau)-\tau)})} \\
			\vspace{0.3cm}\displaystyle{\hspace{1.5cm}-e^{-K(\gamma_r(T+\tau)+\tau)}(M_n^v-\langle x_r(n(\gamma_r(T+\tau)+\tau)),v \rangle)\frac{\psi_0}{N-1}\int_{n(\gamma_r(T+\tau)+\tau)+\tau}^{t}\alpha_{i_1r}(s)e^{-\frac{KN_{i_1}}{N-1}(t-s)}ds}\\
			\vspace{0.3cm}\displaystyle{\hspace{0.8cm}\leq e^{-\frac{KN_{i_1}}{N-1}(t-n(\gamma_r(T+\tau)+\tau)-\tau)}M_n^v +M^v_n(1-e^{-\frac{KN_{i_1}}{N-1}(t-n(\gamma_r(T+\tau)+\tau)-\tau)})}\\
			\vspace{0.3cm}\displaystyle{\hspace{1.5cm}-e^{-K(\gamma_r(T+\tau)+\tau)}(M_n^v-\langle x_r(n(\gamma_r(T+\tau)+\tau)),v \rangle)e^{-K\gamma_r(T+\tau)}\frac{\psi_0}{N-1}\int_{n(\gamma_r(T+\tau)+\tau)+\tau}^{t}\alpha_{i_1r}(s)ds}\\
			\vspace{0.3cm}\displaystyle{\hspace{0.8cm}=M^v_n-e^{-K(2\gamma_r(T+\tau)+\tau)}(M_n^v-\langle x_r(n(\gamma_r(T+\tau)+\tau)),v \rangle)\frac{\psi_0}{N-1}\int_{n(\gamma_r(T+\tau)+\tau)}^{t}\alpha_{i_1r}(s)ds,}
		\end{array}$$
		for all $t\in [n(\gamma_r(T+\tau)+\tau)+\tau,(n+1)(\gamma_r(T+\tau)+\tau)]$. In particular, for $t\in [n(\gamma_r(T+\tau)+\tau)+T+\tau,(n+1)(\gamma_r(T+\tau)+\tau)]$, since the Generalized Persistence Excitation Condition \eqref{PE} implies that 
		$$\int_{n(\gamma_r(T+\tau)+\tau)}^{t}\alpha_{i_1r}(s)ds\geq \int_{n(\gamma_r(T+\tau)+\tau)+\tau}^{n(\gamma_r(T+\tau)+\tau)+T+\tau}\alpha_{i_1r}(s)ds\geq \frac{\tilde{\alpha}}{(1+n(\gamma_r(T+\tau)+\tau)+\tau)^\beta},$$
		we obtain
		\begin{equation}\label{i_1T}
			\langle x_{i_1}(t),v\rangle\leq M^v_n-e^{-K(2\gamma_r(T+\tau)+\tau)}(M_n^v-\langle x_r(n(\gamma_r(T+\tau)+\tau)),v \rangle)\frac{\psi_0\tilde{\alpha}}{(N-1)(1+n(\gamma_r(T+\tau)+\tau)+\tau)^\beta}.
		\end{equation}	
		At this point, if $p=1$, we stop the argument. Otherwise, if $p>1$, using \eqref{i_1T} we estimate, for $t\in [n(\gamma_r(T+\tau)+\tau)+T+2\tau,(n+1)(\gamma_r(T+\tau)+\tau)]$,
		$$\begin{array}{l}
			\vspace{0.3cm}\displaystyle{\frac{d}{d t}  \langle x_{i_2}(t),v\rangle =  \sum_{j \neq i_1, i_2} \chi_{i_2j} b_{i_2j}(t)( \langle x_{j}(t-\tau_{i_2j}(t)),v\rangle-\langle x_{i_2}(t),v\rangle)+ b_{i_2i_1}(t)(\langle x_{i_1}(t-\tau_{i_2i_1}(t)),v\rangle-\langle x_{i_2}(t),v\rangle) }\\
			\vspace{0.3cm}\displaystyle{\hspace{2cm}\leq ( M_{n}^v-\langle x_{i_2}(t),v\rangle)\sum_{j \neq i_1, i_1} \chi_{i_2j} b_{i_2j}(t)} \\
			\displaystyle{+ b_{i_2i_1}(t)\Big (M^v_n-e^{-K(2\gamma_r(T+\tau)+\tau)}(M_n^v-\langle x_r(n(\gamma_r(T+\tau)+\tau)),v \rangle)\frac{\psi_0\tilde{\alpha}}{(N-1)(1+n(\gamma_r(T+\tau)+\tau)+\tau)^\beta}}\\ \vspace{0.3cm}
				\displaystyle{\hspace{16 cm} -\langle x_{i_2}(t),v\rangle\Big).}
		\end{array}$$
		Therefore, arguing as above, 
		$$\begin{array}{l}
			\vspace{0.3cm}\displaystyle{\frac{d}{d t}  \langle x_{i_2}(t),v\rangle \leq \frac{KN_{i_{2}}}{N-1}(M_{n}^v-\langle x_{i_2}(t),v\rangle) -b_{i_2i_1}(t)(M_{n}^v-\langle x_{i_2}(t),v\rangle) + b_{i_2i_1}(t)(M^v_n-\langle x_{i_2}(t),v\rangle)}\\
			\vspace{0.3cm}\displaystyle{\hspace{1.5cm}-b_{i_2i_1}(t)e^{-K(2\gamma_r(T+\tau)+\tau)}(M_n^v-\langle x_r(n(\gamma_r(T+\tau)+\tau)),v \rangle)\frac{\psi_0\tilde{\alpha}}{(N-1)(1+n(\gamma_r(T+\tau)+\tau)+\tau)^\beta}}\\
			\vspace{0.3cm}\displaystyle{\hspace{0.8cm}\leq \frac{KN_{i_{2}}}{N-1}M_{n}^v-\frac{KN_{i_{2}}}{N-1}\langle x_{i_2}(t),v\rangle}\\
            \displaystyle{\hspace{1.5cm}-\alpha_{i_2i_1}(t)e^{-K(2\gamma_r(T+\tau)+\tau)}(M_n^v-\langle x_r(n(\gamma_r(T+\tau)+\tau)),v \rangle)\left(\frac{\psi_0}{N-1}\right)^2\frac{\tilde{\alpha}}{(1+n(\gamma_r(T+\tau)+\tau)+\tau)^\beta}.}
		\end{array}$$
		Again, the Gronwall's inequality yields
		$$ \begin{array}{l}
			\vspace{0.3cm}\displaystyle{\langle x_{i_2}(t),v\rangle\leq M^v_n -e^{-K(2\gamma_r(T+\tau)+{\tau})}(M_n^v-\langle x_r(n(\gamma_r(T+\tau)+\tau)),v \rangle)\left(\frac{\psi_0}{N-1}\right)^2\frac{\tilde{\alpha}}{(1+n(\gamma_r(T+\tau)+\tau)+\tau)^\beta}\times}\\
			\vspace{0.3cm}\displaystyle{\hspace{4cm} \times\int_{n(\gamma_r(T+\tau)+\tau)+T+2\tau}^{t}\alpha_{i_2i_1}(s)e^{-\frac{KN_{i_2}}{N-1}(t-s)}ds}\\
			\vspace{0.3cm}\displaystyle{\hspace{1.5cm}\leq M^v_n -e^{-K(3\gamma_r(T+\tau)-T)}(M_n^v-\langle x_r(n(\gamma_r(T+\tau)+\tau)),v \rangle)\left(\frac{\psi_0}{N-1}\right)^2\frac{\tilde{\alpha}}{(1+n(\gamma_r(T+\tau)+\tau)+\tau)^\beta}\times}\\
			\displaystyle{\hspace{4cm}\times\int_{n(\gamma_r(T+\tau)+\tau)+T+2\tau}^{t}\alpha_{i_2i_1}(s)ds,}
		\end{array}$$
		for all $t\in [n(\gamma_r(T+\tau)+\tau)+T+2\tau,T+2\tau,(n+1)(\gamma_r (T+\tau)+\tau)]$. In particular, for $t\in [n(\gamma_r(T+\tau)+\tau)+2T+2\tau,(n+1)(\gamma_r(T+\tau)+\tau)]$, since from \eqref{PE} it comes that
		$$\int_{n(\gamma_r(T+\tau)+\tau)+T+2\tau}^{t}\alpha_{i_2i_1}(s)ds\geq \int_{n(\gamma_r(T+\tau)+\tau)+T+2\tau}^{n(\gamma_r(T+\tau)+\tau)+2T+2\tau}\alpha_{i_2i_1}(s)ds\geq \frac{\tilde\alpha}{(1+n(\gamma_r(T+\tau)+\tau)+T+2\tau)^\beta},$$
		we can conclude that
		\begin{equation}\label{i_2T}
			\begin{array}{l}
				\vspace{0.3cm}\displaystyle{\langle x_{i_2}(t),v\rangle\leq M^v_n-e^{-K(3\gamma_r (T+\tau)-T)}(M_n^v-\langle x_r(n(\gamma_r(T+\tau)+\tau)),v \rangle)\left(\frac{\psi_0\tilde{\alpha}}{N-1}\right)^2\times}\\
				\vspace{0.3cm}\displaystyle{\hspace{1.8cm}\times\frac{1}{(1+n(\gamma_r(T+\tau)+\tau)+T+2\tau)^\beta}\frac{1}{(1+n(\gamma_r(T+\tau)+\tau)+\tau)^\beta}}\\
				\displaystyle{\hspace{0.8cm}\leq M^v_n-e^{-K(3\gamma_r (T+\tau)-T)}(M_n^v-\langle x_r(n(\gamma_r(T+\tau)+\tau)),v \rangle)\left(\frac{\psi_0\tilde{\alpha}}{(N-1)(1+n(\gamma_r(T+\tau)+\tau)+T+2\tau)^\beta}\right)^2.}
			\end{array}
		\end{equation}
		Finally, iterating the above procedure,
		\begin{equation} \label{5.13onoff}
			\begin{array}{l}
				\vspace{0.3cm}\displaystyle{\langle x_{i_k}(t),v\rangle \leq M^v_{n}-e^{- K((k+1)\gamma_r(T+\tau) -\left(\sum_{l=0}^{k-1}l\right)(T+\tau)+\tau)} (M_n^v-\langle x_r(n(\gamma_r(T+\tau)+\tau)),v \rangle)\times}\\
				\displaystyle{\hspace{2.5cm}\times\left(\frac{\psi_{0}\tilde{\alpha}}{(N-1)(1+n(\gamma_r(T+\tau)+\tau)+(k-1)T+k\tau)^\beta}\right)^{k},}
			\end{array}
		\end{equation}
		for all $1\leq k\leq p$ and for all $t\in [n(\gamma_r (T+\tau)+\tau)+k(T+\tau),(n+1)(\gamma_r (T+\tau)+\tau)]$. In particular, for $k=p$, inequality \eqref{5.13onoff} reads as
        \begin{equation}\label{agenti}
        \begin{split}
           & \langle x_{i}(t),v\rangle \leq M^v_{n}-e^{- K((p+1)\gamma_r(T+\tau) -\left(\sum_{l=0}^{p-1}l\right)(T+\tau)+\tau)} (M_n^v-\langle x_r(n(\gamma_r(T+\tau)+\tau)),v \rangle) \\
           & \hspace{4cm}\left(\frac{\psi_{0}\tilde{\alpha}}{(N-1)(1+n(\gamma_r(T+\tau)+\tau)+(p-1)T+p\tau)^\beta}\right)^{p},
            \end{split}
        \end{equation} 
        for all $t\in [n(\gamma_r (T+\tau)+\tau)+p(T+\tau),(n+1)(\gamma_r (T+\tau)+\tau)]$.
        \\Note that, if the path has length $p=\gamma_r$, using that $\sum_{l=0}^{\gamma_r-1}l=\frac{\gamma_r(\gamma_r-1)}{2}$, estimate \eqref{agenti} gives
		\begin{equation}\label{5.13gammaonoff}
        \begin{array}{l}
        	\vspace{0.3cm}\displaystyle{\langle x_{i}(t),v\rangle \leq M^v_{n}}\\
        	\displaystyle{\hspace{0.1cm}-e^{-K (\frac{1}{2}(\gamma_{r}^2+3\gamma_r)(T+\tau)+\tau)}(M_n^v-\langle x_r(n(\gamma_r(T+\tau)+\tau)),v \rangle)\Big(\frac{\psi_{0}\tilde{\alpha}}{(N-1)(1+n(\gamma_r(T+\tau)+\tau)+(\gamma_r-1)T+\gamma_r\tau)^\beta}\Big)^{\gamma_r},}
        \end{array}
		\end{equation}
		for all $t\in [(n+1)(\gamma_r(T+\tau)+\tau))-\tau,(n+1)(\gamma_r(T+\tau)+\tau))]=I_{n+1}$. Hence, due to the fact that \eqref{5.13gammaonoff} is the coarsest estimate, we conclude that \eqref{5.13gammaonoff} holds for every $i=1,\dots,N$ and $t\in I_{n+1}$. 
\\On the other hand, using analogous arguments one can show that
		\begin{equation*}
		    \begin{array}{l}
		    	\vspace{0.3cm}\displaystyle{\langle x_{i}(t),v\rangle \geq m^v_{n}}\\
		    	\displaystyle{\hspace{0.1cm}+e^{-K(\frac{1}{2}(\gamma_{r}^2+3\gamma_r)(T+\tau)+\tau)}(\langle x_r(n(\gamma_r(T+\tau)+\tau)),v \rangle-m_{n}^v)\Big(\frac{\psi_{0}\tilde{\alpha}}{(N-1)(1+n(\gamma_r(T+\tau)+\tau)+(\gamma_r-1)T+\gamma_r\tau)^\beta}\Big )^{\gamma_r}.}
		    \end{array}
		\end{equation*}
		for all $i=1,\dots,N$ and $t\in I_{n+1}$. Combining this last estimate with \eqref{5.13gammaonoff} we finally obtain \eqref{Bonoff}.
	\end{proof} 

\end{document}
